\documentclass[11pt]{amsart}

\usepackage[T1]{fontenc}
\usepackage{lmodern}
\usepackage{microtype}
\usepackage{mathtools,amssymb,amsthm}
\usepackage[hidelinks]{hyperref}

\hypersetup{
  pdftitle={A stabilizer code admitting no logical Z that meets the hypothesis of Zheng-Liu Proposition 3}
}

\newtheorem{theorem}{Theorem}
\newtheorem{lemma}[theorem]{Lemma}
\newtheorem{proposition}[theorem]{Proposition}
\newtheorem{corollary}[theorem]{Corollary}
\theoremstyle{remark}
\newtheorem*{remark}{Remark}
\theoremstyle{definition}
\newtheorem{definition}[theorem]{Definition}

\newcommand{\PP}{\mathcal P}
\newcommand{\NN}{\mathcal N}
\newcommand{\Zb}{\overline Z}
\newcommand{\Xb}{\overline X}
\newcommand{\Yb}{\overline Y}
\newcommand{\str}[1]{\texttt{#1}}

\title[Failure of the hypothesis of Zheng--Liu Proposition~3]
{A Stabilizer Code Admitting No Logical $Z$\\That Meets the Hypothesis of Zheng--Liu Proposition~3}

\date{}

\subjclass[2020]{81P73, 81P68, 94B05}
\keywords{magic state distillation, stabilizer code, logical Pauli operator, invariant plane,
counterexample}

\begin{document}

\begin{abstract}
Zheng and Liu map a magic-state distillation protocol to a dynamical system on the Bloch
coordinates $(x,y,z)$. Their Proposition~3 shows that the plane $z=0$ is a fixed plane of the map
provided that $s_j\Zb_i$ contains at least one tensor factor $Z$ for every logical operator
$\Zb_i$ and every element $s_j$ of the stabilizer group. Immediately after that proposition they
conjecture, in a hedged sentence of unnumbered running text, that the hypothesis can always be met,
for any stabilizer code, by properly choosing the logical $Z$ operators. For the $[[3,1,1]]$
code $S=\langle \str{XZI},\str{ZXX}\rangle$ used by the same paper, each of the four cosets of $S$
in its normalizer contains a Pauli string with no tensor factor $Z$, so no choice of logical $Z$
meets the condition. Proposition~3 itself is correct and is untouched; the map is never computed
here, and nothing here shows that this code lacks an invariant coordinate plane. What fails is only
the assertion that the hypothesis of Proposition~3 can always be arranged.
\end{abstract}

\maketitle

\section{The claim}\label{sec:claim}

Let $\PP_n$ be the group of $n$-qubit Pauli operators. For a stabilizer code with abelian
stabilizer group $S\le\PP_n$ of rank $n-k$ \cite{G}, write $\NN(S)$ for the normalizer of $S$ in
$\PP_n$ and call the nontrivial cosets of $S$ in $\NN(S)$ the \emph{logical classes} of the code;
there are $4^k$ cosets in all, since $|\NN(S)|=2^{n+k}$ and $|S|=2^{n-k}$ modulo phases.

Zheng and Liu \cite{ZL} send a stabilizer-code magic-state distillation protocol to a polynomial map
of the Bloch vector of the input state --- a viewpoint going back to Rall \cite{R}, on top of the
structural theory of stabilizer-code distillation protocols of Campbell and Browne \cite{CB}. Their
Proposition~3 reads, verbatim:

\begin{quote}
For a stabilizer code $Q$, $z=0$ is a fixed plane in the mapped dynamical system if $s_j\Zb_i$
explicitly contains Pauli $Z$ operators for every logical operator $\Zb_i$ $(i=1,\dots,k)$ and
every stabilizer element $s_j$ $(j=1,\dots,2^{\bar n})$.
\end{quote}

Here $\bar n=n-k$, so the index $j$ runs over the whole group $S$ and not over a generating set.
Immediately after the two-line proof of that proposition, the following sentence appears:

\begin{quote}
This condition is satisfied by most known protocols that we have examined and \emph{we conjecture it
to be generally true for any stabilizer codes by properly choosing the logical $Z$ operators without
modifying the stabilizer groups.}
\end{quote}

It is the universal form of this sentence that is at issue below. Three conventions of \cite{ZL} fix
its content, and we record them because the counterexample is sensitive to all three.

\smallskip\noindent\emph{(i) ``Contains Pauli $Z$'' is about tensor-factor type, and a $Y$ factor
does not count.} The proof of Proposition~1 of \cite{ZL} writes a stabilizer element as
$s_j=(-1)^{\alpha_j}s_j^1\otimes\cdots\otimes s_j^n$ with $s_j^\ell\in\{I,X,Y,Z\}$, and the paper
then counts the weights $w^X,w^Y,w^Z$ of $\Zb_is_j$ separately, $Y$ contributing to the Bloch
monomial $y$ and not to $z$. The factorisation $T^z_i=z\cdot\widetilde T^z_i$ that the proof of
Proposition~3 needs therefore requires a genuine factor $Z$.

\smallskip\noindent\emph{(ii) Phases are irrelevant,} for the same reason: only the multiset of
tensor-factor letters is read. We accordingly work in $\PP_n$ modulo phases throughout, identifying
a Pauli operator with its string over $\{I,X,Y,Z\}$; the product of two strings is the factorwise
product, e.g.\ $\str{XZI}\cdot\str{ZXX}=\str{YYX}$.

\smallskip\noindent\emph{(iii) ``Without modifying the stabilizer groups'' fixes $S$.} A legitimate
choice of logical $Z$ operators is then $k$ elements $\Zb_1,\dots,\Zb_k$ of $\NN(S)$ that pairwise
commute and are independent modulo $S$.

\begin{definition}\label{def:clean}
A Pauli string is \emph{$Z$-free} if none of its tensor factors is $Z$, i.e.\ if it is a string over
$\{I,X,Y\}$. A coset $QS\subseteq\NN(S)$ is
\emph{dirty} if it contains a $Z$-free string, and \emph{clean} otherwise.
\end{definition}

\begin{lemma}[the coset reduction]\label{lem:coset}
Fix $Q\in\NN(S)$. Then $Q$ satisfies the hypothesis of Proposition~3 --- that $s_jQ$ contains a
factor $Z$ for every $s_j\in S$ --- if and only if the coset $QS$ is clean. In particular the
hypothesis depends only on the class of $Q$ in $\NN(S)/S$, and multiplying $Q$ by an element of $S$
buys nothing.
\end{lemma}

\begin{proof}
As $s_j$ runs over $S$ the product $s_jQ$ runs over the coset $QS$ exactly once, because $S$ is a
group. So ``every $s_jQ$ contains a factor $Z$'' says precisely that no member of $QS$ is $Z$-free.
The identity lies in $S$, so $Q$ itself is one of the strings tested.
\end{proof}

By Lemma~\ref{lem:coset}, ``properly choosing the logical $Z$ operators'' is a choice among the
$4^k-1$ nontrivial logical classes: the conjecture asserts that every stabilizer code admits $k$
pairwise commuting, independent-modulo-$S$ elements of $\NN(S)$ whose classes are all clean. We
call such a $k$-tuple a \emph{clean frame}. Note that the trivial class $S$ is always dirty, since
the identity string is $Z$-free; and for a code with no clean class at all, not merely does no valid
frame exist but no single $\Zb_i$ does.

\section{A code of the source paper for which no choice works}\label{sec:witness}

\begin{theorem}\label{thm:311}
Let $S=\langle \str{XZI},\str{ZXX}\rangle$, the stabilizer group of the $[[3,1,1]]$ code worked out
in \cite{ZL}. Then
$S=\{\str{III},\str{XZI},\str{ZXX},\str{YYX}\}$, $|\NN(S)|=16$, and $\NN(S)/S$ has four classes,
each of them dirty:
\[
\begin{array}{llll}
 \text{class of }\str{III}: & \str{III}\cdot\str{III}=\str{III}, &
 \text{class of }\str{IIX}: & \str{IIX}\cdot\str{III}=\str{IIX},\\
 \text{class of }\str{IZZ}: & \str{IZZ}\cdot\str{YYX}=\str{YXY}, &
 \text{class of }\str{IZY}: & \str{IZY}\cdot\str{XZI}=\str{XIY}.
\end{array}
\]
Hence no $Q\in\NN(S)$ satisfies the hypothesis of Proposition~3 for this code, so the universal
claim of the conjectural sentence quoted in Section~\ref{sec:claim} fails.
\end{theorem}

\begin{proof}
The two generators commute --- their symplectic product is
$(1,0,0)\cdot(1,0,0)+(0,1,0)\cdot(0,1,1)=1+1=0$ over $\mathbb F_2$ --- and are independent, so
$S$ is an abelian group of order $4$ with $\str{XZI}\cdot\str{ZXX}=\str{YYX}$, and the code is a
$[[3,1]]$ code, $k=1$. The published logical pair of \cite{ZL} is
$\Zb=\str{IIX}$, $\Xb=\str{IZZ}$; both lie in $\NN(S)$, neither lies in $S$, and they anticommute,
so they are a genuine conjugate pair and the third nontrivial class is that of
$\Yb=\Zb\,\Xb=\str{IZY}$. For $k=1$ these three classes together with $S$ exhaust $\NN(S)/S$, since
$4^k=4$.

Each displayed line is an identity in $\PP_3$ modulo phase whose multiplier is an element of $S$,
and each right-hand side is a string over $\{I,X,Y\}$, i.e.\ $Z$-free. So each of the four cosets
contains a $Z$-free string and is dirty. By Lemma~\ref{lem:coset} no $Q\in\NN(S)$ has $s_jQ$
containing a factor $Z$ for all $s_j\in S$. A~fortiori no clean frame exists. The stabilizer group
has not been modified.
\end{proof}

\section{Scope}\label{sec:scope}

Proposition~3 is a sufficient condition, stated with ``if'' and never with ``iff'', and it is
correct; nothing above touches it. What Theorem~\ref{thm:311} contradicts is only the separate,
unnumbered and hedged sentence asserting that the hypothesis of Proposition~3 can always be
arranged, and in that sentence only the universal quantifier over stabilizer codes. In particular
nothing above says that $z=0$ fails to be a fixed plane of the map for this code: we do not compute
the map at all. Nor is the coordinate-plane route to a fixed plane closed for it, since the analogue
of the condition in the $y$ direction --- every member of the coset carries a factor $Y$ --- is
satisfied by exactly one of its four classes. The other codes and protocols of \cite{ZL} are not
examined here, and nothing here contradicts any numerical result of \cite{ZL}.

The witness code is not ours: \cite{ZL} takes it from Howard and Dawkins~\cite{HD} and reproduces it
verbatim, and what is new here is the statement about it rather than the object.

\section{Reproduction and verification}

Theorem~\ref{thm:311} is a short computation in $\PP_3$ modulo phase: encode a
length-$n$ string as the pair $(a,b)\in\mathbb F_2^n\times\mathbb F_2^n$ with $a_i=1$ iff factor $i$
is $X$ or $Y$ and $b_i=1$ iff it is $Z$ or $Y$; the product is $(a\oplus a',b\oplus b')$, two
strings commute iff $\langle a,b'\rangle+\langle b,a'\rangle=0$ over $\mathbb F_2$, and a string
contains a factor $Z$ iff $\overline a\wedge b\neq0$.

The files \str{verify.py} and \str{verify.output.txt} accompanying this note are a program and a
recorded run of it. From literals only, with no external data, in exact integer arithmetic and with
no randomness, it checks for the code of Theorem~\ref{thm:311} that $S$ is abelian with the printed
element list, that $|\NN(S)|=2^{n+k}$ and $|\NN(S)/S|=4^k$, that each displayed product closes as
Pauli arithmetic with a multiplier genuinely in $S$ and a $Z$-free right-hand side, and that the four
representatives lie in pairwise distinct cosets and so exhaust the classes; independently of those
certificates it then sweeps all of $\NN(S)$ and counts the classes on which the condition holds, in
each of the three coordinate directions. Those are the checks that bear on this paper. The recorded
run also carries checks on other codes and statements, and its labels and commentary refer to
sections, lemmas, propositions and tables of a longer draft rather than to anything printed here;
none of that is claimed in this paper. Nothing in the program computes a dynamical map: what it
decides is the satisfiability of the hypothesis of Proposition~3, not whether $z=0$ is a fixed plane.

\begin{thebibliography}{9}

\bibitem{CB}
E.~T. Campbell and D.~E. Browne,
\emph{On the structure of protocols for magic state distillation},
Theory of Quantum Computation, Communication and Cryptography (TQC 2009),
Lecture Notes in Comput. Sci. \textbf{5906}, Springer, 2009, 20--32.
\href{https://arxiv.org/abs/0908.0838}{arXiv:0908.0838}.

\bibitem{G}
D.~Gottesman,
\emph{Stabilizer codes and quantum error correction},
Ph.D. thesis, California Institute of Technology, 1997.
\href{https://arxiv.org/abs/quant-ph/9705052}{arXiv:quant-ph/9705052}.

\bibitem{HD}
M.~Howard and H.~Dawkins,
\emph{Small codes for magic state distillation},
Eur. Phys. J. D \textbf{70} (2016), 55.
\href{https://doi.org/10.1140/epjd/e2016-60682-y}{doi:10.1140/epjd/e2016-60682-y};
\href{https://arxiv.org/abs/1512.04765}{arXiv:1512.04765}.

\bibitem{R}
P.~Rall,
\emph{Fractal properties of magic state distillation},
preprint, 2017.
\href{https://arxiv.org/abs/1708.09256}{arXiv:1708.09256}.

\bibitem{ZL}
Y.~Zheng and Y.~Liu,
\emph{From magic state distillation to dynamical systems},
Quantum \textbf{9} (2025), 1858.
\href{https://doi.org/10.22331/q-2025-09-15-1858}{doi:10.22331/q-2025-09-15-1858};
\href{https://arxiv.org/abs/2412.04402}{arXiv:2412.04402}. Proposition~3 and the conjectural
sentence following it are quoted from the accepted version, arXiv:2412.04402v2, which carries the
journal acceptance date.

\end{thebibliography}

\end{document}
